% 2025 CUMCM A 问题一首轮审核稿（revision 1）
% 本文件是单问题 LaTeX 片段，不是完整参赛论文；待用户审核后再并入总稿。
\providecommand{\PoneReleasePoint}{(17620,\,0,\,1800)}
\providecommand{\PoneBurstTime}{5.1}
\providecommand{\PoneBurstPoint}{(17188,\,0,\,1736.496)}
\providecommand{\PoneFullStart}{8.056}
\providecommand{\PoneFullEnd}{9.448}
\providecommand{\PoneFullDuration}{1.392}
\providecommand{\PoneCentreStart}{8.013}
\providecommand{\PoneCentreEnd}{9.448}
\providecommand{\PoneCentreDuration}{1.435}
\providecommand{\PoneDurationDifference}{0.043}

\section{问题一模型的建立与求解}

\subsection{问题一题目解读}

以假目标为坐标原点，地面为 $z=0$ 平面。导弹 M1 的初始位置为
$\bm r_{M0}=(20000,0,2000)\,\mathrm m$，并以 $300\,\mathrm{m/s}$ 的速度飞向原点；
无人机 FY1 的初始位置为 $\bm r_{F0}=(17800,0,1800)\,\mathrm m$，按题设以
$120\,\mathrm{m/s}$ 沿等高度直线飞向假目标。真目标是底面圆心为
$(0,200,0)\,\mathrm m$、半径为 $7\,\mathrm m$、高为 $10\,\mathrm m$ 的圆柱体。

FY1 在任务开始后 $1.5\,\mathrm s$ 投放烟幕弹，烟幕弹再经过 $3.6\,\mathrm s$ 起爆，
故绝对起爆时刻为 $5.1\,\mathrm s$。起爆后的烟幕有效区域近似为半径 $10\,\mathrm m$
的球体，球心以 $3\,\mathrm{m/s}$ 竖直下沉，有效期为 $20\,\mathrm s$。因此，问题一可转化为：
在 $[5.1,25.1]$ s 内，求烟幕球能够同时截住 M1 到圆柱真目标所有点视线的时间集合长度。
这一口径比只考察导弹到目标几何中心的单条视线更严格。

\subsection{模型假设}

结合题设给出的短时间运动过程，作如下假设：

\begin{enumerate}
  \item M1 始终沿初始位置到假目标原点的直线，以 $300\,\mathrm{m/s}$ 匀速飞行；
  \item FY1 可瞬时调整到题设航向，随后以 $120\,\mathrm{m/s}$ 等高度匀速直线飞行；
  \item 烟幕弹投放瞬间继承无人机水平速度，初始竖直速度为零，仅受重力作用，取
  $g=9.8\,\mathrm{m/s^2}$；
  \item 忽略空气阻力、风速以及烟幕的水平漂移和形变；
  \item 起爆后烟幕为半径 $10\,\mathrm m$ 的球体，球心以 $3\,\mathrm{m/s}$ 竖直下沉，
  且只在起爆后 $20\,\mathrm s$ 内有效；
  \item 真目标为位置固定的实心圆柱体，半径为 $7\,\mathrm m$，高度为 $10\,\mathrm m$。
\end{enumerate}

\subsection{运动关系与遮蔽判据}

记 $t=0$ 为任务开始时刻。M1 指向原点的单位方向向量为
$-\bm r_{M0}/\lVert\bm r_{M0}\rVert$，故其位置为
\begin{equation}
  \bm r_M(t)=\bm r_{M0}+300t\frac{-\bm r_{M0}}{\lVert\bm r_{M0}\rVert}.
  \label{eq:p1-missile}
\end{equation}
式\eqref{eq:p1-missile}给出了任意时刻的导弹视点。FY1 沿 $x$ 轴负方向飞行，其位置为
\begin{equation}
  \bm r_F(t)=\bm r_{F0}+(-120,0,0)t.
  \label{eq:p1-uav}
\end{equation}
令投放时刻 $t_d=1.5\,\mathrm s$，则由式\eqref{eq:p1-uav}得到投放点
$\bm r_d=\PoneReleasePoint\,\mathrm m$。投放后烟幕弹作平抛运动：
\begin{equation}
  \bm r_B(t)=\bm r_F(t_d)+(-120,0,0)(t-t_d)
  +\left(0,0,-\frac{1}{2}g(t-t_d)^2\right),\qquad t_d\le t\le t_b .
  \label{eq:p1-bomb}
\end{equation}
其中 $t_b=t_d+3.6=\PoneBurstTime\,\mathrm s$。将 $t_b$ 代入式\eqref{eq:p1-bomb}，得起爆点
$\bm r_b=\PoneBurstPoint\,\mathrm m$。因此，烟幕球心在有效期内的位置为
\begin{equation}
  \bm r_C(t)=\bm r_b+\bigl(0,0,-3(t-t_b)\bigr),\qquad t_b\le t\le t_b+20 .
  \label{eq:p1-smoke}
\end{equation}
式\eqref{eq:p1-smoke}同时限定了烟幕的运动方式和有效时间窗。

真目标点集记为
\[
  \mathcal T=\left\{(x,y,z):x^2+(y-200)^2\le 7^2,\ 0\le z\le10\right\}.
\]
对任意目标点 $\bm p\in\mathcal T$，M1 到该点的视线段可写为
$\bm r_M(t)+\lambda[\bm p-\bm r_M(t)]$，其中 $0\le\lambda\le1$。定义烟幕球心到全部视线段的
最不利距离
\begin{equation}
  D(t)=\max_{\bm p\in\mathcal T}\ \min_{0\le\lambda\le1}
  \left\lVert\bm r_C(t)-\left\{\bm r_M(t)+\lambda[\bm p-\bm r_M(t)]\right\}\right\rVert .
  \label{eq:p1-occlusion}
\end{equation}
当且仅当式\eqref{eq:p1-occlusion}满足 $D(t)\le10\,\mathrm m$ 时，半径 $10\,\mathrm m$ 的烟幕球
才对圆柱目标形成完整遮蔽。于是有效遮蔽集合及其时长为
\begin{equation}
  \mathcal I=\left\{t\in[t_b,t_b+20]:D(t)\le10\right\},
  \qquad \Delta t=\operatorname{meas}(\mathcal I).
  \label{eq:p1-duration}
\end{equation}

\subsection{数值求解}

先由式\eqref{eq:p1-uav}和式\eqref{eq:p1-bomb}手算核对，得到投放点
$\PoneReleasePoint\,\mathrm m$、起爆时刻 $\PoneBurstTime\,\mathrm s$ 和起爆点
$\PoneBurstPoint\,\mathrm m$。随后对圆柱侧面及上下底面进行离散，在烟幕有效区间内以
$0.02\,\mathrm s$ 为步长扫描式\eqref{eq:p1-occlusion}的余量 $10-D(t)$；发现符号变化后，
采用 Brent 有界求根法求精进入和离开完整遮蔽状态的时刻，最终按式\eqref{eq:p1-duration}
计算有效区间的长度。将扫描步长依次取为 $0.05\,\mathrm s$、$0.02\,\mathrm s$ 和
$0.01\,\mathrm s$，经边界求精后所得时长的最大差异约为 $1.0\times10^{-12}\,\mathrm s$，
表明时间扫描加密不改变三位小数结果。

圆柱边界依次采用 $180\times17\times17$、$360\times31\times31$ 和
$720\times61\times61$ 三档离散。三档所得完整遮蔽时长分别为
$1.391645\,\mathrm s$、$1.391643\,\mathrm s$ 和 $1.391643\,\mathrm s$；后两档变化小于
$0.001\,\mathrm s$。最细网格下两端点的距离残差均小于 $10^{-7}\,\mathrm m$，且所得区间完全落在
$[5.1,25.1]$ s 内。

作为对照，另取圆柱几何中心 $\bm p_c=(0,200,5)\,\mathrm m$，只计算 M1 到 $\bm p_c$ 的
中心视线是否穿过烟幕球。该对照不作为问题一的主判据。

\subsection{结果与简短分析}

\begin{table}[htbp]
  \centering
  \caption{问题一两种遮蔽口径的计算结果}
  \label{tab:p1-results}
  \label{tab:p1-results-centre}
  \begin{tabular}{lccc}
    \toprule
    判定口径 & 开始时刻/$\mathrm s$ & 结束时刻/$\mathrm s$ & 有效时长/$\mathrm s$ \\
    \midrule
    完整圆柱目标 & \PoneFullStart & \PoneFullEnd & \PoneFullDuration \\
    目标几何中心视线（对照） & \PoneCentreStart & \PoneCentreEnd & \PoneCentreDuration \\
    \bottomrule
  \end{tabular}
\end{table}

由表\ref{tab:p1-results}可知，按完整圆柱目标判据，烟幕在任务开始后
$\PoneFullStart\,\mathrm s$ 开始形成有效遮蔽，在 $\PoneFullEnd\,\mathrm s$ 结束，
故问题一的有效遮蔽时长为 $\boxed{\PoneFullDuration\,\mathrm s}$。若只考察目标几何中心视线，
时长为 $\PoneCentreDuration\,\mathrm s$，比完整目标口径多 $\PoneDurationDifference\,\mathrm s$。
差异主要来自遮蔽初期：烟幕已覆盖目标中心方向，但尚未覆盖圆柱边缘对应的全部视线。因此，
采用中心视线会轻微高估完整遮蔽时间，完整圆柱判据更符合“遮住整个真目标”的要求。

\subsection{针对问题一的摘要素材}

针对给定 FY1 航向、速度及投放—起爆时序，建立导弹、无人机、烟幕弹与烟幕球心的三维运动模型，
并以烟幕球心到“导弹—圆柱目标全部点”视线段集合的最大最短距离构造完整遮蔽判据。通过有效期扫描、
边界求根和圆柱表面加密检验，得到完整目标的有效遮蔽区间为
$[\PoneFullStart,\PoneFullEnd]$ s，持续 $\PoneFullDuration\,\mathrm s$；该结果略短于只考察目标几何中心
视线所得的 $\PoneCentreDuration\,\mathrm s$，说明完整目标判据能够避免中心线近似造成的时长高估。
